\section{Futurists: where the ideas went}
\label{sec:futurists}

\subsection{What the citation record says}
\label{sec:fut-citations}

Figure~\ref{fig:citations} plots the annual citation counts of the four
journal items, taken from OpenAlex on 19 August 2026; absolute counts differ
between bibliographic databases, so only one source is used and only shapes
and orders of magnitude are read from it. The shapes are eloquent. Grubbs'
1969 paper has accumulated about 3\,700 citations and still attracts roughly
170 to 230 a year, more than half a century after publication, a pattern
bibliometricians associate with method classics that outlive their original
field: the citing literature today is dominated not by statistics but by
chemistry, medicine, engineering and environmental science, wherever a
laboratory needs a named, defensible way to discard a measurement. The
applied line is smaller and steeper. Shafer and colleagues stand at roughly
290 citations, a steady trickle from every group that builds a network and
discovers it needs quality assurance; Campbell and colleagues at about 130;
and Leigh and colleagues, youngest of the five, already at about 130 and
cited more in each recent year than either applied predecessor, a
straightforward sign of where the field's active margin now lies.

\begin{figure}[t]
  \centering
  \includegraphics[width=0.98\textwidth]{fig_citations.png}
  \caption{Citations per year for the four journal items on the reading
  list, from OpenAlex, accessed 19 August 2026. Note the different vertical
  scales: the left panel shows a statistics classic in steady state, the
  right panel the applied succession.}
  \label{fig:citations}
\end{figure}

\subsection{From paper to standard to reflex}
\label{sec:fut-standards}

Each reading hardened into infrastructure along its own path. Grubbs' tests
were written into standards, most visibly the ASTM practice on dealing with
outlying observations, and into the reference handbooks and statistical
software through which most users now meet them, usually without reading the
paper. Rosner's generalised ESD travelled the same road as the recommended
repair. The modified z-score of Iglewicz and Hoaglin diffused further and
lost more provenance on the way: it appears in engineering handbooks,
anomaly detection libraries and sensor processing toolkits, frequently with
the 0.6745 misplaced or the 3.5 quoted as if it were a law of nature, and
its origin in a slim quality control booklet is remembered mainly by those
who go looking. Ideas that succeed operationally tend to decay into reflexes;
the constants survive while the caveats fall away, which is exactly how a
batch screening rule ended up running, with inherited thresholds, inside
five minute data streams.

The Mesonet's blueprint became the template for surface network quality
assurance worldwide. Its own authors published a ten year retrospective
extending the battery and the flag taxonomy \citep{fiebrich2010}, syntheses
such as \citet{estevez2011} codified the same test families as recommended
practice for automatic weather stations, and the architecture, automated
graded flags feeding a small human review loop, reappears in national
climate reference networks, agrometeorological services and hydrological
agencies on every continent. Campbell and colleagues' framework was absorbed
even more directly: the observatory generation that followed, most visibly
NEON, built automated plausibility, gap and consistency testing into its
design from the start \citep{taylor2013}, and the paper's vocabulary,
quality assurance versus quality control, provisional versus validated
products, flags carried as data, is now simply how environmental data
producers speak. My own internship system, with its ordered battery of
range, persistence, step, consistency and neighbour checks writing graded
quality codes for later human validation, is a faithful instance of this
lineage, which is precisely why I chose these readings.

\subsection{The current frontier}
\label{sec:fut-frontier}

The decade since Leigh and colleagues splits, on my reading, into three
active fronts. The first is software consolidation. The methods of the 2019
framework and its companion papers shipped as R packages, and the same years
produced PyHydroQC, which pairs detection with model-based correction of
aquatic sensor data \citep{jones2022}, and SaQC, a configuration-driven
quality control engine from the German environmental research community
built so that an agency's entire QC policy is a version-controlled,
executable document \citep{schmidt2023}. This is the quiet but consequential
shift: quality control moving from bespoke scripts toward shared, citable,
configurable infrastructure, which is also the only realistic route by which
threshold provenance, the lineage's weakest joint, becomes auditable.

The second front is machine learning. The anomaly detection literature has
exploded, with deep sequence models, autoencoders and attention
architectures dominating the methods venues; \citet{blazquez2021} give the
panoramic taxonomy. Contact with operational environmental QC remains
thinner than the volume of papers suggests, for reasons the reading list
makes predictable: agencies need per-flag explanations, stable behaviour
under distribution shift, and error rates they can budget, and they possess
few labelled faults to train on, since the whole point of a good QA system
is that confirmed faults are rare. The methods that are crossing over do so
on the lineage's own terms, as additional members of an ensemble whose
verdicts still flow through flags and human review, exactly the pattern
\citet{leigh2019} recommended. I read the current enthusiasm for time series
foundation models the same way: plausibly a better one-step forecaster
inside Equation~\eqref{eq:regression}, not a replacement for the
architecture around it.

The third front is the one my thesis question lives on: making the reference
set itself empirical. Distance-based neighbour selection, inherited from the
Barnes era, is giving way to selection by observed statistical agreement;
fixed plausibility bands are giving way to bands learned from a station's own
validated history; and the field is slowly acquiring what it has always
lacked, public labelled benchmarks on which detectors can be compared the
way Leigh and colleagues compared them privately. Add the unresolved
questions of correction, when a flagged value may be replaced by an estimate
and how that estimate's uncertainty is carried downstream, and of false
alarm budgeting as an explicit design quantity, and the research programme
for the next decade is visible in outline. None of it retires the 1969
paper, the 1993 booklet, the 2000 architecture or the 2013 framework; it
runs on top of them.
